\begin{exercisechunk}
\item \textbf{Barron approximation versus a fixed linear space.}
\label[exercise]{ex:l15-barron-linear-comparison}
\exerciseprofile{3}{2}{3}{\exproof\quad\excalculation\quad\exinterpretation}{%
Prove the orthogonal-family lower-bound mechanism and compare it with
target-dependent ridge approximation.}
Let $e_1,\ldots,e_{2m}$ be orthonormal elements of a Hilbert space $\cH$,
let $V\subseteq\cH$ satisfy $\dim(V)\leq m$, and let $P_V$ be the
orthogonal projection onto $V$.
\begin{enumerate}[(a),beginpenalty=10000]
\item Prove that
\[
 \sum_{j=1}^{2m}\|P_Ve_j\|_{\cH}^2\leq m.
\]
Conclude that some $j\in\{1,\ldots,2m\}$ satisfies
\[
 \inf_{v\in V}\|e_j-v\|_{\cH}\geq\frac1{\sqrt2}.
\]
\item Let $\cF\subseteq\cH$. Suppose that, for every $m\geq1$, there are
$2m$ orthonormal elements $e_1,\ldots,e_{2m}$ such that
$\rho_m e_j\in\cF$ for every $j$, where $\rho_m>0$. Prove that the
Kolmogorov width satisfies
\[
 d_m(\cF;\cH)\geq\frac{\rho_m}{\sqrt2}.
\]
\item Assume that every $f\in\cB_C(Q)$ has an $m$-unit ridge
approximation satisfying
\[
 \|f-f_m\|_2\leq 2C m^{-1/2},
\]
and that every fixed $m$-dimensional linear space $V_m$ satisfies
\[
 \sup_{f\in\cB_C(Q)}
 \inf_{g\in V_m}\|f-g\|_2
 \geq \kappa(C/d)m^{-1/d}.
\]
Give a sufficient width for the ridge approximation to attain error
at most $\epsilon$. Give a necessary dimension for a fixed linear space to
attain worst-case error at most $\epsilon$.
\item Compare the dependence on $\epsilon$ when $d=2$ and when $d=20$.
Which choices in the ridge approximation are allowed to depend on the
target, and which choices in the linear lower bound are fixed?
\item State why the comparison proves neither a smaller minimax regression
risk for a neural-network estimator nor a contradiction with the minimax
lower bound for a full $d$-dimensional smoothness ball. Identify the missing
learning statement in the first case and the different target classes in the
second case.
\end{enumerate}
\begin{hints}
For part (a), expand $P_V$ in an orthonormal basis of $V$ and apply Bessel's
inequality to each basis vector.
\end{hints}
\end{exercisechunk}
